\cleardoublepage
\chapter{TEMPLATE QUERY CYPHER}

Lampiran ini menyajikan template \textit{query} Cypher inti yang digunakan sistem \textit{GraphRAG} Kubernetes untuk mengakses \textit{knowledge graph} di Neo4j, sebagaimana dibahas pada Bab~\ref{chap:implementasi}. Query lengkap beserta seluruh varian (\textit{ablation}, \textit{baseline}) tersedia di berkas \texttt{src/graph/queries.py} pada repositori proyek.

% ─────────────────────────────────────────────────────────────────────────────
\section{Pencarian Exact Match}
% ─────────────────────────────────────────────────────────────────────────────

Query berikut mencari \textit{root node} berdasarkan kecocokan nama persis, digunakan sebelum \textit{vector search} untuk memastikan presisi saat nama \textit{resource} disebutkan secara eksplisit dalam pertanyaan pengguna.

\lstinputlisting[
    language=Cypher,
    caption={Pencarian \textit{Exact Match} berdasarkan Nama \textit{Resource}},
    label={lst:cypher-exact-match}
]{listings/cypher-exact-match.cypher}

% ─────────────────────────────────────────────────────────────────────────────
\section{Traversal Multi-Hop (Schema Dependencies)}
% ─────────────────────────────────────────────────────────────────────────────

Query berikut menelusuri dependensi skema dari \textit{root node} yang ditemukan, menggunakan seluruh 18 tipe \textit{edge} dalam \textit{knowledge graph} untuk membangun konteks yang diberikan ke LLM. Tipe \textit{edge} dan kedalaman maksimum disubstitusi saat \textit{call-time} karena Cypher tidak mendukung parameterisasi \texttt{\$}-\textit{param} untuk nama tipe \textit{relationship} maupun batas \textit{variable-length path}.

\lstinputlisting[
    language=Cypher,
    caption={Traversal Multi-Hop untuk Konteks Skema (\textit{Schema Dependencies})},
    label={lst:cypher-multihop}
]{listings/cypher-multihop-traversal.cypher}

% ─────────────────────────────────────────────────────────────────────────────
\section{Pencarian Vektor (Vector Search)}
% ─────────────────────────────────────────────────────────────────────────────

Query berikut adalah jalur \textit{fallback} saat pencarian \textit{exact match} gagal menemukan \textit{root node}. Pencarian memanfaatkan \textit{Native Vector Index} Neo4j dengan fungsi kemiripan \textit{cosine}, kemudian diperluas dengan traversal multi-hop yang sama.

\lstinputlisting[
    language=Cypher,
    caption={Pencarian Hibrida Vektor dan Graf (\textit{Hybrid Vector + Graph Search})},
    label={lst:cypher-vector-search}
]{listings/cypher-vector-search.cypher}

% ─────────────────────────────────────────────────────────────────────────────
\section{Validasi Required Field}
% ─────────────────────────────────────────────────────────────────────────────

Query berikut mengambil daftar \textit{field} wajib (\textit{required}) suatu \textit{resource} langsung dari \textit{knowledge graph}, digunakan pada Lapisan 3 validasi YAML (\textit{YAMLValidator}).

\lstinputlisting[
    language=Cypher,
    caption={Validasi \textit{Required Field} terhadap \textit{Knowledge Graph}},
    label={lst:cypher-required-fields}
]{listings/cypher-required-fields.cypher}
